\chapter{Incident Investigation, Reporting and Lessons Learned Risk Assessment}
\label{chap:Incident Investigation, Reporting and Lessons Learned Risk Assessment}

%---------------------------- SECTION ----------------------------
\section{Why this assessment matters in construction}

\paragraph{Every incident on a construction site --- from the near miss that goes unreported to the fatality that closes the site for weeks --- carries information. It carries information about the hazards that were present, the controls that failed, the management decisions that created the conditions for failure, and the systemic factors that made those decisions possible. The incident investigation and lessons-learned process is the mechanism through which that information is extracted, analysed, communicated, and used to prevent recurrence. An industry that investigates incidents superficially, that stops its analysis at the immediate cause without following the chain back to root cause, that does not communicate lessons across the organisation, and that does not verify that the lessons have been implemented is condemned to repeat its failures.}

\paragraph{RIDDOR 2013 --- the Reporting of Injuries, Diseases and Dangerous Occurrences Regulations --- is the legal framework through which construction incidents are notified to the Health and Safety Executive and other enforcing authorities. RIDDOR reporting is not merely a legal obligation; it is the mechanism through which the HSE's enforcement programme is activated. A RIDDOR-reported fatality will typically result in a joint investigation by the HSE and the police (under the Work-Related Deaths Protocol); a specified injury will result in an HSE investigation visit within 24 hours; a dangerous occurrence may result in a Prohibition Notice while the investigation is underway. The quality of the construction employer's own internal investigation, and the timeliness and accuracy of the RIDDOR report, directly affects the trajectory of the subsequent regulatory investigation.}

\paragraph{Root cause analysis (RCA) is the discipline that moves incident investigation beyond the immediate ``what happened'' to the deeper ``why it was possible for this to happen.'' The fishbone (Ishikawa) diagram maps the causal factors across categories (Man, Machine, Method, Material, Measurement, Environment) to identify all contributing causes; the 5-Whys technique asks ``why'' repeatedly until the underlying organisational or systemic cause is reached; the bow-tie model maps the event between the threat pathways on the left (hazard to event) and the consequence pathways on the right (event to outcome), identifying where prevention barriers failed and where mitigation barriers were absent. All three techniques are applicable to construction incident investigation and each illuminates different aspects of the causal structure.}

\paragraph{The criminal law dimension of construction incident investigation is significant and often underestimated by construction employers. A construction fatality is simultaneously a RIDDOR-reportable incident, a potential Corporate Manslaughter Act 2007 prosecution, a potential Health and Safety at Work etc.\ Act 1974 prosecution of the employer and of individuals, and a potential inquest. The police and the HSE have joint investigation protocols for work-related deaths; the Crown Prosecution Service and the HSE jointly decide whether to prosecute under the Corporate Manslaughter Act or under health and safety law. Construction employers who conduct their own post-fatality investigation without understanding this legal context --- who instruct staff to cooperate with the police and HSE without legal advice, who allow documents to be taken without a proper understanding of which documents attract legal privilege, and who fail to preserve the scene --- damage their legal position irreparably.}

\barquote{``Every near miss is a gift: a warning that comes without the cost of injury. The organisation that ignores or under-reports its near misses is discarding the only inexpensive lesson the hazard is willing to offer before it presents its invoice in the form of a fatality.''}

\example{Investigation Scenario}{A scaffold operative falls 3.6 metres from a putlog scaffold on a residential refurbishment project, fracturing his pelvis and sustaining a specified injury. The principal contractor's initial investigation identifies ``inadequate edge protection'' as the cause and implements ``improved edge protection checks'' as the corrective action. The HSE investigator who attends the following morning conducts a 5-Whys analysis: Why was there no edge protection? Because the lift was incomplete at the end of the previous day. Why was the lift incomplete? Because the scaffold gang was three boards short and finished early. Why were they three boards short? Because the delivery was short by three boards and this was not escalated. Why was the delivery shortage not escalated? Because the gang foreman did not know the site had a material shortage escalation procedure. Why did he not know the procedure? Because the site induction did not include it and the Construction Phase Plan does not reference it. The corrective action that emerges from the 5-Whys is fundamentally different from ``improved edge protection checks.''

%---------------------------- SECTION ----------------------------
\section{Legal and professional framework}

\paragraph{The incident investigation and reporting legal framework in the UK construction industry draws on five principal instruments: RIDDOR 2013, which mandates what must be reported, to whom, and in what timeframe; the Health and Safety at Work etc.\ Act 1974, which provides the HSE with investigation and enforcement powers; the Corporate Manslaughter and Corporate Homicide Act 2007, which creates a criminal offence for organisations whose gross breach of duty causes death; the Health and Safety (Offences) Act 2008, which expanded the range of offences that can be tried in the Crown Court and increased maximum penalties; and the Health and Safety Offences, Corporate Manslaughter and Food Safety and Hygiene Offences Definitive Guideline 2016 (the Sentencing Guidelines), which provides the framework for calculating fines in health and safety prosecutions based on the organisation's annual turnover.}

\paragraph{The Work-Related Deaths Protocol --- a joint protocol between the HSE, the Chief Fire Officers' Association, the Police, the Crown Prosecution Service, and the Local Government Association --- governs the conduct of investigations into work-related deaths. It establishes who leads the investigation (police for potential corporate manslaughter; HSE for health and safety offences), how evidence is gathered, how the two strands of investigation are coordinated, and how prosecution decisions are made. Construction employers must understand this protocol before a fatality occurs, because the decisions made in the first hours after a fatality --- who is spoken to, what documents are provided, whether the scene is preserved --- are irreversible.}

\begin{itemize}
  \bitem{RIDDOR 2013}{Requires reporting of: deaths arising from workplace accidents (immediately by telephone, then written report within ten days); specified injuries (fractures other than to fingers, thumbs, and toes; amputations; injuries resulting in loss of consciousness; injuries requiring more than 24 hours in hospital; crush injuries to the head or torso; injuries from electric shock; injuries from exposure to substances; injuries requiring resuscitation; loss of consciousness from lack of oxygen); over-seven-day injuries (injuries that prevent the worker from doing their normal duties for more than seven days, excluding the day of the accident --- to be reported within 15 days); dangerous occurrences (a list of 27 specific categories of near-miss events defined in Schedule 2 to RIDDOR, including scaffold collapse, crane overturning, and unintended release of flammable gas); and occupational diseases (vibration white finger, occupational deafness, occupational asthma, contact dermatitis, and others listed in Schedule 3).}
  \bitem{Corporate Manslaughter and Corporate Homicide Act 2007}{Creates a criminal offence for an organisation (including companies, partnerships, and government departments) where the manner in which the organisation's activities are managed or organised causes a person's death, and amounts to a gross breach of a relevant duty of care owed by the organisation to the deceased. The test is whether senior management's conduct fell far below what could reasonably be expected. Conviction results in an unlimited fine, a remedial order, and a publicity order requiring the organisation to publicise the conviction.}
  \bitem{Health and Safety (Offences) Act 2008}{Extended the range of health and safety offences for which individuals can be imprisoned and increased the maximum fines for certain offences. Combined with the 2016 Sentencing Guidelines, the Act significantly increased the personal exposure of directors, managers, and safety advisers who are found to have been complicit in a health and safety failure that causes injury or death.}
  \bitem{Health and Safety Offences Definitive Guideline 2016}{The sentencing guidelines for health and safety offences calculate fines based on the seriousness of the offence (culpability and risk/harm) and the financial means of the defendant. For organisations, fines are calculated as a proportion of annual turnover; for a large organisation with culpability rated as ``High'' and actual harm amounting to a fatality, the starting-point fine may be \pounds4\,million or more. The Guideline has dramatically increased the financial consequences of health and safety convictions since 2016.}
  \bitem{5-Whys technique}{A root cause analysis method in which the investigator asks ``why'' in relation to each causal factor until the underlying organisational or systemic cause is reached. The method is simple, quick, and effective for straightforward incidents; its limitation is that it follows a single causal chain and may miss multiple parallel causal pathways.}
  \bitem{Fishbone (Ishikawa) diagram}{A root cause analysis tool that maps causal factors across multiple categories (typically Man, Machine, Method, Material, Measurement, Environment) radiating from the incident event. The fishbone diagram is better suited than the 5-Whys to complex incidents with multiple causal pathways and is a useful visual communication tool for lessons-learned presentations.}
  \bitem{Bow-tie model}{A risk assessment and incident investigation model that maps the hazard, the top event (the undesired incident), the threat pathways that lead to the top event (on the left, with prevention barriers), and the consequence pathways that lead from the top event to the outcome (on the right, with mitigation barriers). The bow-tie model is particularly useful for identifying which prevention and mitigation barriers failed in a serious incident.}
  \bitem{Near-miss reporting culture}{The ratio of near misses to accidents (expressed in Heinrich's Triangle, though the specific ratios are contested, the principle is broadly accepted) means that every serious injury is preceded by many more near-miss events of the same type. Organisations with high near-miss reporting rates have more data from which to learn; organisations with low near-miss reporting rates have a false sense of safety. Encouraging near-miss reporting requires a no-blame reporting environment in which near misses are treated as learning opportunities, not precursors to disciplinary action.}
\end{itemize}

%---------------------------- SECTION ----------------------------
\section{Conducting the assessment using the five-step method}

\subsection{Step 1 --- Identify Hazards}
\paragraph{Incident investigation and reporting hazards include: failure to recognise a reportable event and therefore failure to notify RIDDOR within the required timeframe; inadequate scene preservation after a serious incident, losing evidence that is essential to the investigation; conducting an internal investigation without legal advice in a way that prejudices the employer's legal position; inadequate root cause analysis that identifies only the immediate cause and implements superficial corrective actions; failure to communicate lessons from one project or site to the wider organisation; and closure of corrective actions on paper without verifying that the changes have been implemented in practice. A specific hazard in the construction context is the temptation to resume work too quickly after a serious incident, before the investigation is complete and the controls verified as adequate, due to programme pressure.}

\subsection{Step 2 --- Decide Who or What Is Affected}
\paragraph{The persons affected by inadequate incident investigation and reporting include: the injured worker, whose future safety on site may depend on the corrective actions identified in the investigation; the workforce as a whole, whose exposure to the hazard that caused the incident continues until the root cause controls are implemented; the employer, who faces regulatory enforcement, criminal prosecution, and civil litigation arising from inadequate investigation and delayed RIDDOR reporting; the HSE inspector who attends the incident scene and whose investigation depends on scene preservation and the accuracy of the employer's initial report; the police, who lead the investigation of any work-related death; and future workers on other sites who would benefit from the lessons learned but will not receive them if the lessons-learned process is not implemented.}

\subsection{Step 3 --- Evaluate Likelihood and Consequence}
\paragraph{The consequence of a failure to report a reportable incident to RIDDOR within the required timeframe is a criminal offence; the HSE has prosecuted construction employers for RIDDOR under-reporting, including the failure to report over-seven-day injuries that were discovered through sickness absence records during a proactive inspection. The consequence of inadequate investigation is the recurrence of a preventable incident; the consequence of recurrence in the case of a serious or fatal hazard is another serious injury or death. The consequence of contaminating the post-fatality scene or making admissions to police without legal advice is the loss of legal defences that could have been preserved.}

\subsection{Step 4 --- Select and Implement Controls}
\paragraph{Controls for incident investigation and reporting must include: a documented incident classification system that identifies which incidents require RIDDOR reporting, immediate telephone notification, or written report only; a scene preservation procedure that is activated immediately after any serious incident and specifies who controls access to the scene and what evidence must be preserved; a legal escalation procedure that requires legal advice to be obtained before any significant document is released to enforcement authorities in a serious incident; a tiered investigation process that applies 5-Whys for minor incidents and fishbone or bow-tie analysis for significant incidents; a corrective action tracking system that assigns ownership, target dates, and evidence-of-completion requirements for each corrective action; and a lessons-learned communication process that distributes investigation outcomes across the organisation in a format that is actionable by site managers and operatives.}

\subsection{Step 5 --- Record and Review}
\paragraph{Investigation records --- scene photographs, witness statements, equipment inspection records, RIDDOR notification copies, investigation reports, and corrective action close-out evidence --- must be maintained in a retrievable format for the limitation period applicable to the potential claims arising from the incident (typically six years for contract claims and three years for personal injury claims from the date of knowledge, but potentially longer). The corrective action register must be reviewed monthly to identify overdue actions; corrective actions that are overdue or have not been verified as implemented must be escalated to the site manager or board, depending on the severity of the underlying hazard.}

\example{Five-Step Application}{Following an over-seven-day injury to a carpenter who cut his thumb severely on a circular saw, a principal contractor implements the five-step investigation process. Step 1 identifies: RIDDOR reporting obligation (over-seven-day); immediate cause (guard not replaced after blade change); root cause through 5-Whys (saw guard removal permitted during blade changes; no procedure for guard replacement before use; site induction covers PPE but not specific saw guard checks). Step 2 identifies: all operatives using circular saws on all company sites. Step 3 rates the risk of future circular saw injury without controls as High (frequent use, High consequence). Step 4 implements: revised site induction to include circular saw guard check; RAMS updated for carpentry to include pre-use saw inspection; toolbox talk delivered on all company sites within five working days; saw guard replacement added to daily pre-use check log. Step 5 specifies: corrective action verified as complete with evidence (updated RAMS, induction records, toolbox talk attendance sheets); monthly review of all open corrective actions.}

%---------------------------- SECTION ----------------------------
\section{Applying the hierarchy of controls}

\paragraph{The hierarchy of controls for incident investigation and reporting is unusual in that the goal is not to prevent the investigation (which is necessary and beneficial) but to ensure that the investigation is conducted thoroughly, that the root causes identified are addressed at the highest possible level of the control hierarchy, and that the corrective actions are verified as effective.}

\begin{itemize}
  \bitem{Elimination}{Eliminate the recurrence of the incident by implementing elimination-level controls where they are identified as feasible: remove the hazardous substance, machine, or process that caused the incident from the site or from the method of work. An investigation that leads to a corrective action at PPE level, when the root cause analysis reveals that the hazard could be eliminated, has failed to apply the hierarchy.}
  \bitem{Substitution}{Implement substitution-level corrective actions where elimination is not feasible: substitute the machine, process, or method that caused the incident with a less hazardous alternative. An investigation that correctly identifies the need for substitution but leaves the corrective action as ``additional training'' is not applying the hierarchy.}
  \bitem{Engineering controls}{Implement engineering-level corrective actions for hazards identified in the investigation that are addressable by physical means: guarding, interlocking, edge protection, safe access systems. Engineering controls should be the default corrective action for any hazard that can be addressed physically.}
  \bitem{Administrative controls}{Implement administrative corrective actions that address the systemic and organisational root causes identified in the investigation: revised RAMS, updated induction, new procedure, toolbox talk, revised supervision arrangement. Administrative controls must be monitored for compliance and must not be the only corrective action for hazards addressable at a higher level.}
  \bitem{PPE}{Implement PPE requirements identified in the investigation where they are appropriate and proportionate; PPE is a valid and important last line of defence, but the investigation must ensure that PPE is not used as the only corrective action for a hazard that requires a higher-order control. A corrective action that says ``ensure operatives wear gloves'' when the root cause is that the machine lacks adequate guarding is not an adequate investigation outcome.}
\end{itemize}

%---------------------------- SECTION ----------------------------
\section{Cadence of assessment and review triggers}

\paragraph{Incident investigation and reporting has a reactive cadence driven by events rather than a fixed calendar cycle; however, the review of investigation quality, corrective action completion, and lessons-learned dissemination must be embedded in the regular management cycle.}

\begin{itemize}
  \bitem{Immediately after any incident}{Scene preservation must be activated immediately, RIDDOR notification assessed, and a preliminary investigation commenced within 24 hours. The on-call manager must be notified of any serious injury or dangerous occurrence regardless of the time of day.}
  \bitem{Within 24 hours for specified injuries and dangerous occurrences}{A competent investigator must attend the scene within 24 hours of any specified injury or dangerous occurrence to conduct a preliminary investigation, gather witness accounts while memories are fresh, and assess the RIDDOR reporting obligation.}
  \bitem{Within 15 days for over-seven-day injuries}{RIDDOR requires over-seven-day injuries to be reported within 15 days; the investigation should be completed in time to inform the RIDDOR report.}
  \bitem{Monthly corrective action review}{All open corrective actions from incident investigations must be reviewed monthly to identify overdue items and escalate them for completion.}
  \bitem{Quarterly lessons-learned review}{Incident investigation outcomes and lessons learned from across the company's projects should be collated and reviewed quarterly, with a summary disseminated to all site managers and safety advisers.}
  \bitem{Annual review of investigation and reporting procedures}{The incident investigation procedure, RIDDOR classification guide, and lessons-learned communication process should be reviewed annually to ensure they remain current with changes in legislation, guidance, and the company's risk profile.}
\end{itemize}

%---------------------------- SECTION ----------------------------
\section{Common failure modes}

\begin{itemize}
  \bitem{RIDDOR not reported or reported late}{The most common RIDDOR failure in construction is the failure to report over-seven-day injuries, often because the employer does not track the number of days the worker is absent and does not trigger the RIDDOR notification when the seven-day threshold is crossed. Over-seven-day injuries that appear on sickness absence records but not in the RIDDOR report are a consistent finding of HSE proactive inspection visits.}
  \bitem{Investigation stops at immediate cause}{The most common investigation failure is the identification of the immediate cause (``operative was not wearing a harness'') without pursuing the question of why the immediate cause occurred (why was the operative not wearing a harness? because it was not available; why was it not available? because it was not included in the equipment order; why was it not included? because the RAMS did not specify it; why did the RAMS not specify it? because the RAMS was generic and not reviewed for this specific task). Stopping at the immediate cause produces corrective actions that address the symptom rather than the disease.}
  \bitem{Corrective actions not verified as implemented}{Corrective actions are assigned on paper, delegated to site managers, and never followed up. The investigation report is filed; the corrective actions are assumed to be complete; the hazard that caused the original incident remains uncontrolled. The recurrence investigation identifies the same root causes as the first.}
  \bitem{Lessons not shared across the organisation}{Investigation reports that are filed on the site server and not shared with other projects allow the same incident to recur on a different site. An organisation with 20 active construction sites that does not have a process for sharing lessons learned across all sites will learn the same lesson 20 times.}
  \bitem{Scene disturbed before investigation is complete}{The pressure to resume work after a serious incident frequently leads to the disturbance or clearance of the scene before the investigation is complete. Evidence that could have established the precise cause of the incident (the position of a guard, the state of a scaffold fitting, the reading on a gas detector) is irretrievably lost.}
\end{itemize}

\example{Failure Pattern}{A scaffolder sustains a fatal fall from a scaffold structure that collapses. The HSE and police attend jointly under the Work-Related Deaths Protocol. During the HSE investigation it emerges that: a near-miss report filed three weeks earlier identified a loose standard base plate on the same scaffold; the near-miss was investigated by the site manager who noted ``operative to retighten''; no competent scaffold inspection was conducted following the near-miss; the corrective action was not verified as completed. The scaffold inspector's log shows the last inspection was six days before the fatality. The Corporate Manslaughter Act 2007 investigation by the Crown Prosecution Service reviews the management system: the near-miss reporting system existed; the near-miss was reported; the investigation was inadequate; the corrective action was not completed; the management oversight of scaffold inspection compliance was insufficient. The company is indicted under the Corporate Manslaughter and Corporate Homicide Act 2007.}

%---------------------------- CHAPTER SUMMARY ----------------------------
\newpage
\section{Chapter Summary}

\begin{table}[H]
\centering
\begin{tabularx}{\textwidth}{>{\raggedright\arraybackslash}p{3.2cm}X}
\toprule
\textbf{Assessment Element} & \textbf{Operational Requirement} \\
\midrule
Legal/Professional Basis & RIDDOR 2013; Health and Safety at Work etc.\ Act 1974; Corporate Manslaughter and Corporate Homicide Act 2007; Health and Safety (Offences) Act 2008; Health and Safety Offences Definitive Guideline 2016; Work-Related Deaths Protocol. \\
Method & Tiered investigation process (5-Whys for minor incidents; fishbone or bow-tie for significant incidents); RIDDOR classification system; root cause analysis; corrective action tracking; lessons-learned dissemination. \\
Controls & Scene preservation procedure; RIDDOR notification timetable and classification guide; legal escalation procedure for serious incidents; root cause analysis training for investigators; corrective action register with ownership and evidence requirements; lessons-learned communication process. \\
Cadence & Immediately for scene preservation; within 24 hours for specified injuries; within 15 days for over-seven-day injuries; monthly corrective action review; quarterly lessons-learned review; annual procedure review. \\
Assurance & RIDDOR compliance audit; corrective action close-out verification; lessons-learned dissemination records; investigation quality audit by competent external reviewer. \\
\bottomrule
\end{tabularx}
\end{table}

\begin{itemize}
  \bitem{Key takeaway 1}{RIDDOR reporting is a legal obligation with specific timeframes; over-seven-day injuries are the most commonly missed category and must be tracked through the sickness absence management system as well as the incident reporting system.}
  \bitem{Key takeaway 2}{Root cause analysis must go beyond the immediate cause to the organisational and systemic causes of the incident; an investigation that stops at ``failure to wear PPE'' has not identified the management and systemic failures that made PPE omission possible.}
  \bitem{Key takeaway 3}{Scene preservation is a critical first response action after any serious incident; evidence disturbed before the investigation is complete cannot be recovered, and its loss may affect the employer's legal position.}
  \bitem{Key takeaway 4}{Corrective actions must be verified as implemented in practice, not merely marked as complete on the investigation report; a corrective action that exists on paper but has not changed the physical reality of the site provides no risk reduction.}
  \bitem{Key takeaway 5}{Lessons learned from incident investigations must be systematically shared across all sites and projects; an organisation that investigates incidents thoroughly but does not share the outcomes has invested in the investigation and discarded the return.}
\end{itemize}

\barquote{In Chapter~\ref{chap:Building a Risk-Mature Construction Business --- Audit, Cadence and Continuous Improvement}, we bring together the themes of the book to examine how construction organisations build and sustain a culture of risk maturity through audit, cadence, and continuous improvement.}
